\paragraph{Du codage ternaire au drapeau marqué.}
\label{inc-ampl-bandera}
Les régions fournissent également les codes $w_\pi,w_e,w_\varphi$ et
$w_{\rm APP}$ dans le repère hérité. L'application
\[
 \gamma_{A_W}:w\longmapsto(w,wA_W)
\]
conserve le mot de six composantes comme première coordonnée et ajoute
sa transformation par $A_W$. La relation $A_W^2=-I_6$ fixe la
compatibilité quadratique de cette réalisation. Le pas suivant,
$c\mapsto\operatorname{supp}(c)$, conserve les positions non nulles.
Pour les mots de poids six, le support identifie exactement la paire
$\{c,-c\}$, comme le prouve la distance minimale du code. Ainsi,
le quotient par le signe produit les hexades, tandis que le code complet
reste disponible pour les opérations qui requièrent ses symboles.

Dans cette coordinatisation, les deux constructions de la tétrade coïncident :
\begin{equation}
 \{i:K_i\geq729\}
 =H_e\cap P_\pi=B_K=\{4,6,9,10\}.
 \label{inc-ampl-cuaterna}
\end{equation}
Le membre de gauche provient d'une comparaison entière des composantes
de $K$ ; celui de droite, de l'intersection des supports codés de
propagation et de clôture. Le seuil appartient à la coordinatisation par blocs
déclarée et reste explicité dans la définition du lecteur. Le
registre intégral conserve les valeurs sur lesquelles il agit, même lorsque
des seuils distincts produisent le même sous-ensemble. L'égalité exprime
donc une compatibilité entre des applications spécifiées.

Le support $H^-_\alpha$ complète cette tétrade en une hexade. La
condition d'appartenance au support $P_\pi$, avec les tailles
d'intersection $(4,3,2)$ relativement à $H_e,H_\varphi,H_{\rm APP}$,
sélectionne une unique hexade. Les preuves suivantes réunissent les deux
déterminations : celle obtenue par le signe de $Z_0$ et celle obtenue
par l'incidence. Le résultat s'organise sous la forme
\begin{equation}
 \widetilde{\mathcal I}_*
 =\bigl(B_K\subset H^-_\alpha;
         P_\pi,H_e,H_\varphi,H_{\rm APP}\bigr).
 \label{inc-ampl-marco}
\end{equation}
L'inclusion est le drapeau ; les quatre supports supplémentaires constituent
son repère marqué. Les permutations simultanées transportent le repère et
conservent ses inclusions et ses intersections. L'information pertinente réside
dans cette configuration transportée et dans sa classe d'isomorphisme.

Les cinq régions de $\pi$ retrouvent ici leur différenciation interne.
Elles partagent $w_\pi$ comme code ternaire visible, mais conservent des émissions,
des signatures et des multiplicités distinctes. Dans la réalisation binaire sur une
dodécade marquée, la région transversale correspond à l'octade dont
l'intersection avec la dodécade est la tétrade ; les quatre autres régions
correspondent aux relèvements des quatre hexades de son étoile. La
région négative est distinguée par $H^-_\alpha$ et l'ordre de coordinatisation
individualise les trois régions positives. Cette correspondance conserve
les rôles régionaux et la distribution $432+4\cdot144=1008$, tout en
réalisant la partition d'incidence $1+4$.
